% Разработчики: сотрудники кафедры программного обеспечения вычислительной техники и автоматизированных систем БГТУ им. В.Г. Шухова
%
% Данный файл содержит требования к содержанию и оформлению дипломных работ и магистерских диссертаций кафедры ПОВТАС БГТУ им. В.Г. Шухова


% !TeX program = lualatex

\documentclass[14pt]{extarticle}

\usepackage[drawstamps=true]{bstunote}

\usepackage[a4paper, left=23mm, right=9mm, top=9mm, bottom=24mm]{geometry}

\usepackage{iftex}

\ifTUTeX\else % Проверяем: не используем ли мы PDFLaTeX
\usepackage[T2A]{fontenc}	% Используем кириллицу в PDFLaTeX
\usepackage[utf8]{inputenc} % Кодировка исходного файла в PDFLaTeX - UTF-8
\fi

\usepackage[english, russian]{babel} % Подключаем многоязычную локализацию; russian = правила переноса, надписей (например, "Содержание")

% Шрифт в рамках
\ifTUTeX % Если используется компилятор XeLaTeX или LuaLaTeX
\usepackage{fontspec}
\newfontface{\GOSTFontFamily}{GOST-type-A} [ % Название шрифта, который будет в рамках
Path = ./fonts/,
Extension = .ttf,
FakeSlant=0.3 	% Используем искусственный наклон
]
\StampFont{\GOSTFontFamily}
\else
% В PDFLaTeX для рамок лучше подобрать шрифт, похожий на чертёжный GOST Type A. Но такого там нет, поэтому выбираем стандартный шрифт без засечек.
\StampFont{\fontencoding{T2A}\fontfamily{cmss}\fontshape{it}\selectfont}
\fi

\ifTUTeX % Если используется компилятор XeLaTeX или LuaLaTeX

\defaultfontfeatures{Mapping=tex-text}

% Шрифт с засечками
\setmainfont{PT-Astra-Serif}[
Path = ./fonts/PT-Astra-Serif/,
Extension = .ttf,
UprightFont = *-Regular,
BoldFont = *-Bold,
ItalicFont = *-Italic,
BoldItalicFont = *-BoldItalic]

% Шрифт без засечек
\setsansfont{PT-Astra-Sans}[
Path = ./fonts/PT-Astra-Sans/,
Extension = .ttf,
UprightFont = *-Regular,
BoldFont = *-Bold,
ItalicFont = *-Italic,
BoldItalicFont = *-BoldItalic]

% Моноширинный шрифт
\setmonofont{PT-Mono}[
Path = ./fonts/PT-Mono/,
Extension = .ttf,
BoldFont = *-Bold,
UprightFont = *-Regular,
Scale=0.9
]

%\setmainfont{Times New Roman} % Шрифт с засечками
%\setsansfont{Arial}	% Шрифт без засечек
%\setmonofont[Scale=0.9]{Consolas} % Моноширинный шрифт

%\setmainfont{PT Astra Serif} % Шрифт с засечками
%\setsansfont{PT Astra Sans}	% Шрифт без засечек
%\setmonofont[Scale=0.9]{PT Mono} % Моноширинный шрифт
\fi

\ifLuaTeX % Если используется компилятор LuaLaTeX
% В LuaLaTeX включаем продвинутую микротипографику 
\usepackage[protrusion=true, expansion=true]{microtype}
\fi

\usepackage{float}

% Параметры оформления списка литературы при использовании пакета biblatex
%\usepackage[parentracker=true,
%	backend=biber,
%	style=gost-numeric, % стиль по ГОСТ Р 7.0.100–2018
%	language=auto,
%	autolang=other,
%	sorting=none % Порядок цитирования в тексте
%]{biblatex}
%\addbibresource{bibliography.bib} % Файл с литературными источниками
%\setlength{\biblabelsep}{3pt} % Расстояние после номера в списке литературы

\usepackage[colorlinks=true,linkcolor=blue,urlcolor=blue,citecolor=green!70!black,hypertexnames=false,breaklinks=true]{hyperref}
\usepackage{longtable}
\usepackage{tcolorbox}

% Настройка цветных рамок. Эти рамки нужны только в методичке. В пояснительной записке их делать не нужно.
\tcbset{%
	colback=green!10, 			% Цвет фона
	colframe=green!50!black, 	% Цвет рамки
	fonttitle=\bfseries, 		% Жирный шрифт для заголовка
	coltitle=green!10,    		% Цвет заголовка
	before upper={\setlength{\parindent}{0.75cm}} % Абзацный отступ
}

\usepackage{tikz}

\usetikzlibrary{positioning}
\usetikzlibrary{shapes.geometric}
\usetikzlibrary{shapes.misc}
\usetikzlibrary{calc}
\usetikzlibrary{chains}
\usetikzlibrary{matrix}
\usetikzlibrary{decorations.text}


\StampAuthor{Пахомов В.А.}{3.06.26\,} % "\," - пробел в конце нужен только для наклонного шрифта 
\StampCheckBy{Осипов О.В.}{3.06.26\,}
%\StampConsultant{Картамышев С.В.}{3.06.26\,}
\StampNormControl{Бондаренко Т.В.}{11.06.26\,}
% Если ФИО не помещается в графу, то подбираем шрифт внутри команд \fontsize{15pt}{12pt}\selectfont
% Пример:
%\StampNormControl{\fontsize{9.5pt}{6pt}\selectfont Бондаренко\kern0.1emТ\kern-0.2em.В\kern-0.1em.}{11.06.26\,}
\StampApprovedBy{Поляков В.М.}{25.06.26\,}

\StampWorkTheme{Разработка программного обеспе\-чения для прогнозирования погоды на основе анализа спутниковых снимков с использованием методов машинного обучения}

% Если название ВКР не вмещается в графу, то подбираем шрифт и междустрочное расстояние командой \fontsize. Ещё можно добавить пе-ре-но-сы по сло-гам командой "\-". Но лучше обойтись без переносов по слогам, если это возможно.
% Пример:
%\StampWorkTheme{\fontsize{11pt}{10pt}\selectfont Разработка программного обеспечения для прог\-нозирования погоды на основе анализа спутниковых снимков с ис\-пользованием методов машинного обучения}
\StampOrganization{БГТУ~им.~В.Г.~Шухова, \makebox{ПВ-223}}

\begin{document}
%%% -------------------------------------------------------------------------------------
%%% ------------------------------- Титульный лист --------------------------------------
%%% -------------------------------------------------------------------------------------
\newgeometry{left=3cm, right=1.5cm, top=1.5cm, bottom=1.5cm}
\newpage
\begin{titlepage}
\begin{center}
	\small
	\textbf{МИНОБРНАУКИ РОССИИ} \\[0.5em]
	ФЕДЕРАЛЬНОЕ ГОСУДАРСТВЕННОЕ БЮДЖЕТНОЕ ОБРАЗОВАТЕЛЬНОЕ УЧРЕЖДЕНИЕ 
	ВЫСШЕГО ОБРАЗОВАНИЯ \\[0.3em]
	\textbf{<<БЕЛГОРОДСКИЙ ГОСУДАРСТВЕННЫЙ ТЕХНОЛОГИЧЕСКИЙ\\УНИВЕРСИТЕТ~им.~В.Г.~ШУХОВА>>}\\
	\textbf{(БГТУ~им.~В.Г.~Шухова)}
\end{center}
\vspace{1em}
\begin{list}{}{\setlength{\leftmargin}{0cm}\setlength{\parsep}{1pt}}
	\small\sloppy
	\item\textbf{Институт} информационных технологий и управляющих систем
	\item\textbf{Кафедра} программного обеспечения вычислительной техники и автоматизированных \makebox{систем}
	\item\textbf{Направление подготовки} 09.03.04 Программная инженерия
	\item\textbf{Направленность (профиль) образовательной программы} Разработка программно-информационных \makebox{систем}
\end{list}
\vspace{0.5cm}	
\begin{center}
\textbf{ВЫПУСКНАЯ КВАЛИФИКАЦИОННАЯ РАБОТА} \\[0em]
на тему: \\[1em]
\textbf{Разработка сервиса хранения аудиоданных с применением алгоритмов упаковки и потоковой распаковки} 
\end{center}
\vfill
\begin{flushleft}
	\small
	\begin{tabular}{@{}l l@{}}
		\hspace{0cm}\textbf{Студент:}         & Пахомов Владислав Андреевич \\
		\hspace{0cm}\textbf{Зав. кафедрой:}   & канд.\,техн.\ наук, доц. Поляков Владимир Михайлович \\
		\hspace{0cm}\textbf{Руководитель:}    & канд.\,физ.-мат.\ наук, доц. Осипов Олег Васильевич
	\end{tabular}
\end{flushleft}
\vspace{1em}
{\small
\hspace*{10em}\textbf{К защите допустить} \\[-0.5em]
\hspace*{10em}\textbf{Зав. кафедрой \underline{\hspace{4cm}} /Поляков В.М./} \\[0.3em]
\hspace*{12em}\textbf{<<\underline{\hspace{1cm}}>>\underline{\hspace{3cm}} \the\year \ г.}
}

\vspace{1em}
{\centering\small\textbf{Белгород \the\year \ г.}\par}
\end{titlepage}

%%% -------------------------------------------------------------------------------------
%%% ------------------------------- Листы с заданием ------------------------------------
%%% -------------------------------------------------------------------------------------
\begin{center}
	\small
	\textbf{МИНОБРНАУКИ РОССИИ} \\[0.5em]
	ФЕДЕРАЛЬНОЕ ГОСУДАРСТВЕННОЕ БЮДЖЕТНОЕ ОБРАЗОВАТЕЛЬНОЕ УЧРЕЖДЕНИЕ 
	ВЫСШЕГО ОБРАЗОВАНИЯ \\[0.3em]
	\textbf{<<БЕЛГОРОДСКИЙ ГОСУДАРСТВЕННЫЙ ТЕХНОЛОГИЧЕСКИЙ\\УНИВЕРСИТЕТ~им.~В.Г.~ШУХОВА>>}\\
	\textbf{(БГТУ~им.~В.Г.~Шухова)}
\end{center}%
\begin{list}{}{\setlength{\leftmargin}{0cm}\setlength{\parsep}{1pt}}
	\small\sloppy
	\item\textbf{Институт} информационных технологий и управляющих систем
	\item\textbf{Кафедра} программного обеспечения вычислительной техники и автоматизированных \makebox{систем}
	\item\textbf{Направление подготовки} 09.03.04 Программная инженерия
	\item\textbf{Направленность (профиль) образовательной программы} Разработка программно-информационных систем
\end{list}
\vspace{-0.5cm}
{\small\noindent
\hspace*{8.0cm}Утверждаю:\\
\hspace*{8cm}Зав. кафедрой \hrulefill \,
\hspace*{8.0cm}<<\underline{\hspace{1cm}}>> \underline{\hspace{3cm}} \the\year \ г.
}
\begin{center}
	\small
	\textbf{ЗАДАНИЕ} \\
	на выпускную квалификационную работу студента \\[0.3em]
	\noindent
	\underline{\makebox[\linewidth][c]{\textit{Пахомова Владислава Андреевича}}}
\end{center}
\vspace{1em}
\begin{enumerate}[label=\arabic*., leftmargin=0pt, labelsep=0.25em, itemindent=13pt, itemsep=6pt] % Настройка нумерации
	\small\linespread{1}\selectfont
	\item Вид выпускной квалификационной работы (ВКР): \textit {бакалаврская работа.}
	\item Тема ВКР \textit{Разработка программного обеспечения для прогнозирования погоды на основе анализа спутниковых снимков с использованием методов машинного обучения} 
	утверждена приказом по университету от <<28>> января \the\year \ г. \textnumero 2/124.
	\item Срок сдачи студентом законченной ВКР: <<12>> июня \the\year \ г.
	\item Исходные данные: ключевые слова, название алгоритмов и методов, технологии разработки и т.п. Пример:  рекомендательные системы, подходы построения рекомендательных систем, методы матричной факторизации, оптимизационные методы обучения модели, аффинные преобразования, трёхмерная графика, машинное обучение, клиент-серверное приложение, генеративная нейронная сеть.
	\item Содержание ВКР:
	Здесь должны быть перечислены названия разделов (глав) пояснительной записки. Пример: \par
	\hspace{13pt}Введение,\par
	\hspace{13pt}Описание предметной области прогнозирования погоды,\par
	\hspace{13pt}Разработка архитектуры программного обеспечения прогнозирования погоды,\par
	\hspace{13pt}Разработка алгоритмов и программного обеспечения прогнозирования погоды,\par
	\hspace{13pt}Тестирование программной системы прогнозирования погоды,\par
	\hspace{13pt}Руководство пользователя разработанной системы прогнозирования погоды,\par
	\hspace{13pt}Заключение,\par
	\hspace{13pt}Список литературы,\par
	\hspace{13pt}Приложение А,\par
	\hspace{13pt}Приложение Б.\par
	\newpage
	\item Перечень графического материала: 1278 рисунков, 1785 таблиц. Презентация включает: постановку задачи, актуальность, технологии, численные эксперименты, внешний вид программы, заключение.
\end{enumerate}
\advance\hoffset-1.5cm % Сдвигаем всю страницу влево на 1.5 см
{\small\noindent Консультанты по работе с указанием относящихся к ним разделов:\\%
% Если консультантов нет, то эту таблицу и строку над ней нужно удалить
\begin{tabular}{|p{5cm}|p{4cm}|p{2.9cm}|p{2.9cm}|}
	\hline
	\parbox[c][2cm][c]{5cm}{\centering\textbf{Раздел}} &
	\parbox[c][2em][c]{4cm}{\centering\textbf{Консультант}} &
	\parbox[c][2em][c]{2.9cm}{\centering\textbf{Задание выдал} (подпись, дата)} &
	\parbox[c][2em][c]{2.9cm}{\centering\textbf{Задание принял} (подпись, дата)} \\
	\hline
	&  &  &  \\
	\hline
	%Введение & Сидоров А.И. &  &  \\
	%\hline
\end{tabular}
}

\vspace{1em}
\noindent {\small Дата выдачи задания: <<17>> \textit{января} \the\year\ г. }\\
\vspace{1em}
\noindent
\hspace*{9cm} 
\begin{minipage}{6cm} 
	\centering
	\underline{\hspace{\linewidth}} \\[-0.3em]
	{\footnotesize (подпись руководителя)}
\end{minipage}

\vspace{1em}

\noindent {\small Задание принял (приняла) к исполнению }\\
\vspace{1em}
\noindent
\hspace*{9cm}
\begin{minipage}{6cm}
	\centering
	\underline{\hspace{\linewidth}} \\[-0.3em]
	{\footnotesize (подпись студента)}
\end{minipage}
\vspace{1em}

{\centering \bf \small КАЛЕНДАРНЫЙ ПЛАН \par}%
{\small\linespread{1}\selectfont%
\begin{longtable}{|@{}c@{}|p{8.3cm}|c|c|}%
\hline%
{\parbox[c][1.7cm][c]{0.8cm}{\centering \bf \textnumero{} \\ п/п}} &
\parbox{8.3cm}{\centering \bf Наименование этапов работы} &
\parbox{3.4cm}{\centering \bf  Срок выполнения этапов работы} &
{\centering \bf  Примечание} \\
\endfirsthead
\hline%
{\parbox[c][1.7cm][c]{0.8cm}{\centering \bf \textnumero{} \\ п/п}} &
\parbox{8.3cm}{\centering \bf Наименование этапов работы} &
\parbox{3.4cm}{\centering \bf  Срок выполнения этапов работы} &
{\centering \bf  Примечание} \\
\endhead
	\hline
	1 & Анализ предметной области. Постановка задачи разработки системы прогнозирования погоды методами машинного обучения. & 18.01.26 -- 15.02.26 & Выполнено \\
	\hline
	2 & Разработка алгоритмов, структур нейросетей и архитектуры системы прогнозирования погоды. & 15.02.26 -- 20.03.26 & Выполнено \\
	\hline
	3 & Разработка программной части системы прогнозирования погоды. & 20.03.26 -- 29.03.26 & Выполнено \\
	\hline
	4 & Разработка клиентской части системы прогнозирования погоды. & 29.03.26 -- 29.04.26 & Выполнено \\
	\hline
	5 & Тестирование программного обеспечения для прогнозирования погоды. & 29.04.26 -- 17.05.26 & Выполнено \\
	\hline
	6 & Оформление пояснительной записки. Подготовка выступления. & 17.05.26 -- 10.06.26 & Выполнено \\
	\hline
	5 & Тестирование программного обеспечения для прогнозирования погоды. & 29.04.26 -- 17.05.26 & Выполнено \\
	\hline
	6 & Оформление пояснительной записки. Подготовка выступления. & 17.05.26 -- 10.06.26 & Выполнено \\
	
	\hline
\end{longtable}
}
\vspace{1em}
{\small
\noindent\makebox[0.6\textwidth]{Студент (ка) \hrulefill} \textit{Пахомов Владислав Андреевич}} \\[-0.5em]
\noindent
\hspace*{5cm}{\footnotesize (подпись)}
\hspace{5cm}{\footnotesize (Ф.И.О.)}

\vspace{2em}

{\small
	\noindent\makebox[0.6\textwidth]{Руководитель ВКР \hrulefill} \textit{Осипов Олег Васильевич}}\\[-0.5em]
\noindent
\hspace*{5cm}{\footnotesize (подпись)}
\hspace{5cm}{\footnotesize (Ф.И.О.)}


%%% -------------------------------------------------------------------------------------
%%% ----------------------------- Лист проверки заимствований ---------------------------
%%% -------------------------------------------------------------------------------------
\newpage
\advance\hoffset 1.5cm % Сдвигаем всю страницу обратно на 1.5 см вправо
\begin{center}
	\textbf{Результаты проверки ЭВ ВКР на заимствование}
\end{center}

\vspace{1em}

\noindent Кафедра программного обеспечения вычислительной техники и автоматизированных систем\\[0.5em]

\noindent
Студент (ка)\,\makebox[0cm][l]{\underline{\hspace{13.4cm}}}%
\parbox[t]{3.5cm}{\centering \makebox{Пахомов В.А.}\\[-0.2em] \footnotesize (Ф.И.О.)\par}%
\parbox[t]{4.5cm}{\centering \,\\[-0.2em] \footnotesize (подпись)\par}%
\parbox[t]{2.2cm}{\centering ПВ-223\\[-0.2em] \footnotesize (группа)\par}%
\parbox[t]{2.5cm}{\centering 16.06.\the\year\\[-0.2em] \footnotesize (дата)\par}%


\vspace{1em}

\noindent
Тема ВКР: Разработка сервиса хранения аудиоданных с применением алгоритмов упаковки и потоковой распаковки.

\vspace{1em}

\noindent
ВКР прошла проверку на объём заимствований.\\
Итоговая оценка заимствований: \textbf{0\%}.

\vspace{3em}
\noindent Работу проверил\\
\noindent
\makebox[0.0cm][l]{\underline{\hspace{\linewidth}}}%
\parbox[t]{5.4cm}{\makebox[5.4cm][l]{канд. физ.-мат. наук, доцент}\\[-0.2em]%
{\centering\footnotesize (уч. степень, должность)\par}}%
\parbox[t]{4.0cm}{\centering \,\\[-0.2em] \footnotesize (подпись)\par}%
\parbox[t]{2.5cm}{\centering 16.06.\the\year\\[-0.2em] \footnotesize (дата) \par}%
\parbox[t]{4.5cm}{\centering \makebox{Хлопов А.М.}\\[-0.2em] \footnotesize (Ф.И.О.)\par}%

\vspace{3em}
\noindent Руководитель ВКР\\
\noindent
\makebox[0.0cm][l]{\underline{\hspace{\linewidth}}}%
\parbox[t]{5.4cm}{\makebox[5.4cm][l]{канд. социол. наук, доцент}\\[-0.2em]%
{\centering\footnotesize (уч. степень, должность)\par}}%
\parbox[t]{4.0cm}{\, \\[-0.2em]\centering\footnotesize (подпись)\par}%
\parbox[t]{2.5cm}{\centering 16.06.\the\year \\[-0.2em] \footnotesize (дата)\par}%
\parbox[t]{4.5cm}{\centering\makebox{Островский А.М.}\\[-0.2em]\centering\footnotesize (Ф.И.О.)\par}%

\restoregeometry

\clearpage
\initpagenumbering
\setcounter{page}{4} % Текущая страница - четвёртая

\section*{OПРЕДЕЛЕНИЯ, СOКРAЩЕНИЯ И OБOЗНAЧЕНИЯ}
\par РФ --- Российская Федерация.
\par СВО --- специальная военная операция.
\par ВУЗ --- высшее учебное заведение.
\par ВКР --- выпускная квалификационная работа.
\par ИТ --- информационные технологии.
\par БД --- база данных.
\par ПО --- программное обеспечение.
\par СУБД --- система управления базами данных.
\par ООП --- объектно-ориентированное программирование.
\par ЭВМ --- электронно-вычислительная машина.
\par КОП --- код операции.
\par АЛУ --- арифметико-логическое устройство.
\par ПЛИС --- программируемая логическая интегральная схема.
\par БПФ --- быстрое преобразование Фурье.
\par СЛАУ --- система линейных алгебраических уравнений.
\par ОС --- операционная система.
\par API (Application Programming Interface) --- программный интерфейс, обеспечивающий взаимодействие между компонентами программного обеспечения.
\par RGB (red, green, blue) --- цветовой формат машинной графики, основанный на смешении компонентов красного, зелёного и синего цветов.
\par WWW (World Wide Web) --- всемирная паутина, система взаимосвязанных гипертекстовых документов.
\par PDF (Portable Document Format) --- переносимый формат электронных документов.
\par HTML (HyperText Markup Language) --- язык гипертекстовой разметки, который используется для создания и структурирования веб-страниц.
\par SQL (Structured Query Language) --- язык структурированных запросов, используемый для работы с базами данных.
\par VPN (Virtual Private Network) --- виртуальная частная сеть, обеспечивающая защищённое соединение через интернет.
\par IoT (Internet of Things) --- интернет вещей, концепция подключения различных устройств к интернету для обмена данными.
\par JSON (JavaScript Object Notation) --- текстовый формат обмена данными, который используется для хранения и передачи структурированной информации между системами и приложениями.


% Содержание
\renewcommand{\contentsname}{\rucontentsname}
\tableofcontents

\section*[ВВЕДЕНИЕ]{ВВЕДЕНИЕ} 

Мир современного человека всё больше переходит в цифровое пространство, и это не удивительно.
Цифровое пространство обладает неоспоримыми преимуществами при работе с информацией:

\begin{itemize}[label=$\bullet$]
	\item Дешёвое хранение информации, а значит её большой объём
	\item Большая скорость передачи информации
	\item Возможность обрабатывать большое количество информации очень быстро
	\item Качество хранимой информации
\end{itemize}

Эти факторы так или иначе и стали решающими
в решении хранения данных как в быту человека, так и в бизнес-процессах и в творчестве.

Самым ярким и приближённым к цифровому пространству примером является профессия программиста.
На это влияет как сама специфика и предназначение профессии, так и вышеописанные процессы.
Можно ли себе представить программиста, который совсем не владеет компьютером? 
Конечно же нет. Возьмут ли на работу программиста, который продолжает хранить весь свой 
код на листах бумаги? Только если эта компания почему-то любит запускать очень старые ракеты. 
Современного конкурентноспособного разработчика сложно представить без умных текстовых редакторов, 
хранения кода в облаке, автоматизации процессов доставки клиентам, нейронных сетей помогающих 
в написании кода (а иногда ещё и пишущих за них).

Сама профессия по личному мнению автора лежит на стыке творческого процесса и бизнес процессов 
в разной степени важности в зависимости от разрабатываемого продукта.
Она требует как изобретательности, креативности и нестандартных подходов при решении различных задач, 
так и соответствия конкретным рамкам, условиям и отлаженным процессам. 

На примере профессии программиста можно увидеть, насколько сильно может интегрироваться цифровое 
пространство в сферу требующую как творческого подхода, так и точности,
и насколько много пользы оно может принести при решении задач. Важно, однако, отметить, что
цифровое пространство всё ещё не может полностью вытеснить другие подходы. Это такой же инструмент, 
как, например, молоток. Оно решает конкретные проблемы, связанные преимущественно с информацией.

Одна из сфер, которая очень похожа на программирование, является музыкальная сфера. 
Она также требует творческого подхода, но также задаёт в себе определённые рамки.

И да, за последнее время интеграция цифрового пространства в этой сфере тоже сильно продвинулась. 
Сервисы музыкального стриминга, согласно исследованиям, используют уже более трети 
россиян, а рост аудитории продолжается. В процессе написания музыки всё чаще используются 
DAW, умеющие работать с продвинутым оборудованием записи; синтез звука и его обработка тоже проходили 
очень долгий путь эволюции. 

Однако на рынке продюссирования музыки почему-то очень мало инструментов, которые помогают 
в разработке аудиоданных. Множество продюссеров всё ещё складывают свои проекты в крайне неорганизованное
хранилище часто представляющее собой обычную папку на локальных персональных компьютерах. А 
если данные будут утеряны, то обычно с этим ничего поделать нельзя и их приходится восстанавливать
"по памяти".

И такие проблемы очень легко решить. Тем более, для них уже есть решение. И их можно 
позаимствовать из той же ранее расмотренной профессии программиста. Облачное хранилище данных
позволит избежать проблемы утери данных а также позволит делиться ими касанием пальца, что 
тоже повысит эффективность общения с коллегами. Хранение предыдущих версий позволит 
возвращаться к старым концепциям, пересматривать их, рефлексировать и адаптировать в разработке.
Распространение данных в централизованной платформе гораздо более удобно для конечного потребителя
в лице дистрибьютора или слушателя. Дополнение платформы в виде бизнес-сущностей позволит 
дистрибьюторам при публикации полученных композиций.

\newcounter{lastPageNumber}
\setcounter{lastPageNumber}{\value{page}} % Сохраняем номер последней значащей страницы

% СЛУЖЕБНАЯ ИНФОРМАЦИЯ ДЛЯ ПРОВЕРКИ НА НОРМКОНТРОЛЬ
\newpage
\setcounter{page}{1} % Начинаем подсчёт страниц с единицы
% Эта информация (списки картинок, таблиц, листингов) нужна только для прохождения нормоконтроля.
% При компиляции итогового варианта ВКР для печати эти страницы нужно удалить!
\renewcommand{\ruappendixname}{СЛУЖЕБНАЯ ИНФОРМАЦИЯ}
\listoffigures
{\color{red!60!black} \bf
Список картинок нужен только для прохождения нормоконтроля. \par 
В диплом он не вшивается!	
}

\listoftables
{\color{red!60!black} \bf
	Список таблиц нужен только для прохождения нормоконтроля. \par
	В диплом он не вшивается!	
}

\lstlistoflistings
{\color{red!60!black} \bf
	Список листингов нужен только для прохождения нормоконтроля. \par
	В диплом он не вшивается!	
}
\setcounter{page}{\value{lastPageNumber}} % Восстанавливаем номер последней значащей страницы, потому что LaTeX будет подставлять в графу "Листов" окончательное значение счётчика page.

\end{document}
